% 题证摘录；来源与整理边界见相同编号分题文件。

\paragraph{4.12 Representations of $S_n$.}
In this subsection we give a description of the representations of the symmetric group $S_n$ for any $n$.
\begin{definition}[4.35]
A partition $\lambda$ of $n$ is a representation of $n$ in the form $n=\lambda_1+\lambda_2+\cdots+\lambda_p$, where $\lambda_i$ are positive integers, and $\lambda_i\geq\lambda_{i+1}$.
\end{definition}
To such $\lambda$ we will attach a Young diagram $Y_\lambda$, which is the union of rectangles $-i\leq y\leq-i+1$, $0\leq x\leq\lambda_i$ in the coordinate plane, for $i=1,\ldots,p$. Clearly, $Y_\lambda$ is a collection of $n$ unit squares. A Young tableau corresponding to $Y_\lambda$ is the result of filling the numbers $1,\ldots,n$ into the squares of $Y_\lambda$ in some way (without repetitions). For example, we will consider the Young tableau $T_\lambda$ obtained by filling in the numbers in the increasing order, left to right, top to bottom.

We can define two subgroups of $S_n$ corresponding to $T_\lambda$:
\begin{enumerate}
\item The row subgroup $P_\lambda$: the subgroup which maps every element of $\{1,\ldots,n\}$ into an element standing in the same row in $T_\lambda$.
\item The column subgroup $Q_\lambda$: the subgroup which maps every element of $\{1,\ldots,n\}$ into an element standing in the same column in $T_\lambda$.
\end{enumerate}
Clearly, $P_\lambda\cap Q_\lambda=\{1\}$. Define the Young projectors:
\[
a_\lambda:=\frac1{|P_\lambda|}\sum_{g\in P_\lambda}g,\qquad
b_\lambda:=\frac1{|Q_\lambda|}\sum_{g\in Q_\lambda}(-1)^g g,
\]
where $(-1)^g$ denotes the sign of the permutation $g$. Set $c_\lambda=a_\lambda b_\lambda$. Since $P_\lambda\cap Q_\lambda=\{1\}$, this element is nonzero.

The irreducible representations of $S_n$ are described by the following theorem.
\begin{theorem}[4.36]
The subspace $V_\lambda:=\C[S_n]c_\lambda$ of $\C[S_n]$ is an irreducible representation of $S_n$ under left multiplication. Every irreducible representation of $S_n$ is isomorphic to $V_\lambda$ for a unique $\lambda$.
\end{theorem}
The modules $V_\lambda$ are called the Specht modules.
\paragraph{4.13 Proof of Theorem 4.36.}
\begin{lemma}[4.40]
Let $x\in\C[S_n]$. Then $a_\lambda xb_\lambda=\ell_\lambda(x)c_\lambda$, where $\ell_\lambda$ is a linear function.
\end{lemma}
\begin{proof}
If $g\in P_\lambda Q_\lambda$, then $g$ has a unique representation as $pq$, $p\in P_\lambda$, $q\in Q_\lambda$, so $a_\lambda gb_\lambda=(-1)^q c_\lambda$. Thus, to prove the required statement, we need to show that if $g$ is a permutation which is not in $P_\lambda Q_\lambda$ then $a_\lambda gb_\lambda=0$.

To show this, it is sufficient to find a transposition $t$ such that $t\in P_\lambda$ and $g^{-1}tg\in Q_\lambda$; then
\[
a_\lambda gb_\lambda=a_\lambda tgb_\lambda
=a_\lambda g(g^{-1}tg)b_\lambda=-a_\lambda gb_\lambda,
\]
so $a_\lambda gb_\lambda=0$. In other words, we have to find two elements $i,j$ standing in the same row in the tableau $T=T_\lambda$, and in the same column in the tableau $T'=gT$ (where $gT$ is the tableau of the same shape as $T$ obtained by permuting the entries of $T$ by the permutation $g$). Thus, it suffices to show that if such a pair does not exist, then $g\in P_\lambda Q_\lambda$, i.e., there exists $p\in P_\lambda$, $q'\in Q'_\lambda:=gQ_\lambda g^{-1}$ such that $pT=q'T'$ (so that $g=pq^{-1}$, $q=g^{-1}q'g\in Q_\lambda$).

Any two elements in the first row of $T$ must be in different columns of $T'$, so there exists $q'_1\in Q'_\lambda$ which moves all these elements to the first row. So there is $p_1\in P_\lambda$ such that $p_1T$ and $q'_1T'$ have the same first row. Now do the same procedure with the second row, finding elements $p_2,q'_2$ such that $p_2p_1T$ and $q'_2q'_1T'$ have the same first two rows. Continuing so, we will construct the desired elements $p,q'$. The lemma is proved.
\end{proof}
Let us introduce the lexicographic ordering on partitions: $\lambda>\mu$ if the first nonvanishing $\lambda_i-\mu_i$ is positive.
\begin{lemma}[4.41]
If $\lambda>\mu$ then $a_\lambda\C[S_n]b_\mu=0$.
\end{lemma}
\begin{proof}
Similarly to the previous lemma, it suffices to show that for any $g\in S_n$ there exists a transposition $t\in P_\lambda$ such that $g^{-1}tg\in Q_\mu$. Let $T=T_\lambda$ and $T'=gT_\mu$. We claim that there are two integers which are in the same row of $T$ and the same column of $T'$. Indeed, if $\lambda_1>\mu_1$, this is clear by the pigeonhole principle (already for the first row). Otherwise, if $\lambda_1=\mu_1$, like in the proof of the previous lemma, we can find elements $p_1\in P_\lambda$, $q'_1\in gQ_\mu g^{-1}$ such that $p_1T$ and $q'_1T'$ have the same first row, and repeat the argument for the second row, and so on. Eventually, having done $i-1$ such steps, we'll have $\lambda_i>\mu_i$, which means that some two elements of the $i$-th row of the first tableau are in the same column of the second tableau, completing the proof.
\end{proof}
\begin{lemma}[4.42]
$c_\lambda$ is proportional to an idempotent. Namely, $c_\lambda^2=\frac{n!}{\dim V_\lambda}c_\lambda$.
\end{lemma}
\begin{proof}
Lemma 4.40 implies that $c_\lambda^2$ is proportional to $c_\lambda$. Also, it is easy to see that the trace of $c_\lambda$ in the regular representation is $n!$ (as the coefficient of the identity element in $c_\lambda$ is $1$). This implies the statement.
\end{proof}
\begin{lemma}[4.43]
Let $A$ be an algebra and $e$ be an idempotent in $A$. Then for any left $A$-module $M$, one has $\Hom_A(Ae,M)\cong eM$ (namely, $x\in eM$ corresponds to $f_x:Ae\to M$ given by $f_x(a)=ax$, $a\in Ae$).
\end{lemma}
\begin{proof}
Note that $1-e$ is also an idempotent in $A$. Thus the statement immediately follows from the fact that $\Hom_A(A,M)\cong M$ and the decomposition $A=Ae\oplus A(1-e)$.
\end{proof}
\begin{proof}[Completion of Theorem 4.36]
Now we are ready to prove Theorem 4.36. Let $\lambda\geq\mu$. Then by Lemmas 4.42, 4.43
\[
\Hom_{S_n}(V_\lambda,V_\mu)
=\Hom_{S_n}(\C[S_n]c_\lambda,\C[S_n]c_\mu)
=c_\lambda\C[S_n]c_\mu.
\]
The latter space is zero for $\lambda>\mu$ by Lemma 4.41, and $1$-dimensional if $\lambda=\mu$ by Lemmas 4.40 and 4.42. Therefore, $V_\lambda$ are irreducible, and $V_\lambda$ is not isomorphic to $V_\mu$ if $\lambda\ne\mu$. Since the number of partitions equals the number of conjugacy classes in $S_n$, the representations $V_\lambda$ exhaust all the irreducible representations of $S_n$. The theorem is proved.
\end{proof}
